\input amstex
\documentstyle{amsppt}
\magnification 1200
\pageno 49
\define\ogf{{G(F)}}
\define\ozf{Z_G(F)}
\define\oglf{{GL(2,E)}}
\define\oa{{A_G(F)}}
\define\oslf{{H(E)}}
\define\slf{{H(F)}}
\define\eps{\epsilon}
\define\okk{{\Cal O}_K}
\define\off{{\Cal O}_F}
\define\offp{{\Cal O}_{F'}}
\define\qqq{\Bbb Q}
\define\ppp{\Bbb P}
\define\ccc{\Bbb C}
\define\rrr{\Bbb R}
\define\aaa{\Bbb A}
\define\nnn{\Bbb N}
\define\unt{\underline{T}}
\define\unx{\underline{x}}
\define\oga{{G(\aaa )}}
\define\oba{{B_G(\aaa )}}
\define\oza{{Z_G(\aaa )}}
\define\ogla{{GL(2,\aaa_L )}}
\define\osla{{H(\aaa_L )}}
\define\sla{H(\aaa )}
\define\oee{{\Cal O}_{E}}
\define\oeep{{\Cal O}_{E'}}
\define\VV{\Cal V}
\define\HH{\Cal H}
\hsize 6.5 true in
\vsize 8.7 true in
\vskip .3in
\subheading{\S 8 On the 3-regular germ for $G_2$}
\vskip .1in
Let $G=G_2$ throughtout this section.

The only essential
base divisor which can contribute to the 3-regular
germ for $G_2$ via Igusa theory is $E_{\emptyset,\beta,\zeta}$,
which has $\beta(E_{\emptyset,\beta,\zeta})=4$.
It must be $F$-rational since $\sigma(E_{\emptyset,\beta,\zeta})$
is also a $3$-regular base divisor, by the observation (I) of \S 4
and the case $G=G_2$ of conjecture 3.4 (which we proved in \S 3
above).
Recall that the action of the Galois
group $\Gamma$ on $T_G$ and on $Z$ factors through a finite group
$\Gamma_T$.
To satisfy the hypotheses (5.1-4), one may use
the calculations of \S 2 to ``average'' the coordinates of
$Z_{\emptyset,\beta,\zeta}$
(or a suitable ${\Bbb Z}$-linear combination of them)
over $\Gamma_T$ to get $F$-rational coordinates on this divisor.
\par
The branch $E_{\emptyset,\beta,\zeta}$ is defined by
$\lambda =0$ on the regular subset $Z_{\emptyset,\beta,\zeta}$:
$\{x(\zeta)\not= 0,w(\zeta)\not= 0,z(\alpha)\not= 0\}\subset Z$.
If we use the local coordinates
$z(\alpha)$, $z(\beta)$, $w(\gamma)$, $w(\delta)$,
$w(\eps)$, $w(\zeta)$, and $x(\zeta)$ on
$Z_{\emptyset,\beta,\zeta}$ to compute the form
$\omega_{E}$ on $E=E_{\emptyset,\beta,\zeta}$, we get
$$
\split
\omega_{E_{\emptyset,\beta,\zeta}}
&={\omega_{{Z}_1}\over \lambda^\beta {dz(\beta)\over z(\beta)}}\\
&={{d\lambda \wedge dx(\alpha) \wedge ...\wedge dx(\zeta)}
\over \lambda^4 {dz(\beta)\over z(\beta)}}\\
&= {{x(\zeta)^2w(\zeta)^4}\over w(\delta)^7z(\alpha)^2}
dz(\alpha)\wedge dw(\gamma)\wedge ...\wedge dw(\zeta)\wedge dx(\zeta).
\endsplit
\tag 8.1
$$
The divisor $\lambda =0$ occurs on
$Z_{\emptyset,\beta,\zeta}$ as $z(\beta)=0$ with multiplicity 2.
On $E_{\emptyset,\beta,\zeta}$, $\lambda =z(\beta)=x(\epsilon)=x(\delta)=
x(\gamma)=x(\beta)=x(\alpha)=0$. 

By \S\S 4-5 and formula 0.6, the contribution of the 
3-regular germ to the $\kappa$-orbital
integral is therefore, 
$\kappa(inv(\gamma_G^0,\gamma_{G,\lambda}))D_H(\gamma_{H,\lambda})^{-1}$ times
$$
|\lambda|^3\sum_{\theta}\theta(\lambda)
F_1(\theta,4,\pi^*(f)\kappa_\theta(t_*))=
|\lambda|^3\sum_{\theta}\theta(\lambda)
\int_{E_{\emptyset,\beta,\zeta}(F)}\pi^*(f)\kappa_\theta(t_*)
|\omega_{E_{\emptyset,\beta,\zeta}}|.
$$
Note that only $\theta =1$ contributes to the sum by \cite{H1, ch 6, \S 2}.
Recall that, on $Z_{\emptyset,\beta,\zeta}$,
$$
|\omega_{E_{\emptyset,\beta,\zeta}}|
={|x(\zeta)|^6|z(\alpha)|^{10}\over{|w(\delta)|^7}}
dz(\alpha)dw(\gamma)dw(\delta)dw(\eps)dw(\zeta)dx(\zeta)
dv(\alpha)...dv(\zeta).
$$
We may assume that the support of $f$ meets ${\Cal O}_3^{st}$ and no other
unipotent orbit. In this case, by Formula 0.6(b), we have
$$
\int_{E_{\emptyset,\beta,\zeta}(F)}\pi^*(f)\kappa_\theta(t_*)
|\omega_{E_{\emptyset,\beta,\zeta}}|
=\int_{{\Cal O}_3(F)}f
\int_{E_{\emptyset,\beta,\zeta}(u)}\kappa_1(t_*)
|\omega_{E_{\emptyset,\beta,\zeta}(u)}||\omega_{{\Cal O}_3}|,
\tag 8.2
$$
which is a multiple of the $\kappa$-stable unipotent
orbital integral of $f$ associated to the $3-$regular orbit
of $G_2$. It is shown in the appendix that the $3$-regular orbit is stable.
\proclaim{Theorem 8.3} Suppose $T_G$ is a non-split Cartan subgroup
of $G_2$ which splits over a
Galois extension for which the Galois group action on $T_G$ 
is realized by a cyclic subgroup of $\Omega$ of order three. 
Assume that $f\in C_c^\infty(G(F))$ satisfies $supp(f)\cap {\Cal O}_3
\not= \emptyset$ and the support of $f$ does not meet other unipotent orbit
of $G(F)$. Then there is an $f^H$, smooth on
$H(F)$, for which, for all regular $\gamma_G\in T_G(F)$
in a sufficiently small neighborhood of $1$, we have 
$$
\Delta_{G/H}^{T_G,\kappa}(\gamma_H,\gamma_G)
\Phi^{T_G,\kappa}(\gamma_G,f)=
\Phi^{T_H,st}(\gamma_H,f^H),
$$
where $\kappa$, $T_G$ satisfy $H_{T_G,\kappa}=PGL(3)$. 
(In particular, $\kappa$ is a cubic character and $T_G$ is a cyclic
``cubic'' elliptic Cartan.)
\endproclaim
This will be proven below. 
\vskip .2in
{\bf Proof of Theorem 8.3}
\vskip .1in
Note that by \cite{Ro}, the germ on the
right-hand side, $\Gamma_{{\Cal O}_H^{st}}^{T_H,st}(\gamma_H)$, is a 
non-zero constant. Therefore, the theorem will follow if we can show that,
as a function of the tangent vector $X$ of $C$,
the two sides are equal up to a constant.

Let $u=\theta_{\beta}(1)\in G_2$.
One may explicitly compute the orbit $ad({\frak g}_2)(X_\beta)$.
From this, we conclude that $u^{G_2}$
is 6-dimensional. Therefore ${\Cal O}_3^{st}=u^{G_2}$. 
There is only one $F$-orbit of $u$ contained in the stable orbit of $u$
(as is verified in the appendix), so the $3$-regular orbit
of $G_2$ is stable.
\par
To get the local coordinates on the fiber 
$E_{\emptyset,\beta,\zeta}(u)\subset E_{\emptyset,\beta,\zeta}$ over
$u$ in the 3-regular orbit ${\Cal O}_3$ take the local
coordinates in the ``right'' divisor $E_{\emptyset,\beta,\zeta}$, 
then divide $\omega_{E_{\emptyset,\beta,\zeta}}$ by a Haar measure
$\omega_{{\Cal O}_3}$ on ${\Cal O}_3$.
Analogously to Hales' lemma 6.4,
the subvariety $\prod_{\sigma\in \Sigma_G-\Sigma_{H_1}}U_\sigma$
of $G_2$ yields local coordinates for the $3$-regular orbit,
where $\Sigma_{H_1}^+=\{\beta,\epsilon,\zeta\}$.
We use the form on ${\Cal O}_3$,
$$
\split
|\omega_{{\Cal O}_3}|
&:={|\bigwedge_{\sigma\in \Sigma_{{\Cal O}_3}}dx(\sigma)\wedge dv(\sigma)|
\over |\lambda|^2}\\
&={|x(\zeta)|\over |w(\zeta)||z(\alpha)|}
dz(\alpha)dw(\gamma)dw(\delta)dv(\alpha)dv(\gamma)dv(\delta),
\endsplit
\tag 8.4
$$
where $\Sigma_{{\Cal O}_3}=\Sigma_{G}^+ -\Sigma_{H_1}^+$.
On the generic fiber $E_{\emptyset,\beta,\zeta}(u)$, $u\in {\Cal O}_3$,
$$
|\omega_{E_{\emptyset,\beta,\zeta}(u)}|=
{{|x(\zeta)||w(\zeta)|^5}\over{|z(\alpha)||w(\delta)|^7}}
dw(\epsilon)dw(\zeta)dx(\zeta)dv(\beta)dv(\epsilon)dv(\zeta). 
\tag 8.5
$$
We need to rewrite this last form as an expression
$$
\omega_{E_{\emptyset,\beta,\zeta}(u)}=
|p_{{\Cal O}_3,T_G}(X)|\omega'_{E_{\emptyset,\beta,\zeta}(u)},
$$
where $p_{{\Cal O}_3,T_G}(X)$ is a homogeneous polynomial in the coordinates
of tangent vector $X$ to be computed below
and $\omega'_{E_{\emptyset,\beta,\zeta}(u)}$
is independent of the tangent direction. 
\par
By an easy calculation given in the appendix,
our assumption on $T_G$ implies that $\kappa$ is a character
of order $3$.
\par
Once we do this it will follow that
$$
\split
&D_H(\gamma_{\lambda,H})\sum_{{\Cal O}\subset {\Cal O}_3^{st}}
\mu_{\Cal O}(f)\Gamma_{\Cal O}^{T_G,\kappa}(\gamma_\lambda)
=|\lambda|^2\kappa(inv(\gamma_G^0,\gamma_\lambda))
\int_{E_{\emptyset,\beta,\zeta}}
\pi^*(f)\kappa_1(t_*)|\omega_{{\Cal E}_{\emptyset,\beta,\zeta}(F)}|\\
&=|\lambda|^2\kappa(inv(\gamma_G^0,\gamma_\lambda))
|p_{{\Cal O}_3,T_G}(X)|
\sum_{{\Cal O}\subset {\Cal O}_3^{st}}\int_{{\Cal O}_3(F)}f
\int_{E_{\emptyset,\beta,\zeta}(u)(F)}\kappa(t_*)
|\omega'_{E_{\emptyset,\beta,\zeta}(u)}||\omega_{\Cal O}|,
\endsplit
\tag 8.6
$$
where we normalize $\omega_{{\Cal O}}$ consistent with the
Langlands-Shelstad transfer factor as in \cite{H2}. 

To determine the general form of
$p_{{\Cal O}_3,T_G}(X)$ we use the coordinate relations in (1.22).
We find that $p_{{\Cal O}_3,T_G}(X)$ is a product of the form
$$
p_{{\Cal O}_3,T_G}(X)=\prod_{\sigma \in\Sigma_H^+}\sigma(X)^{e_\sigma},
\tag 8.7
$$
where $\sum_{\sigma \in\Sigma_H^+}e_\sigma=3$.
One interesting thing to note is that the product runs over roots
``in $H$''. 
\par
To complete the matching, it remains only to 
verify the local translation invariance property in the hypothesis
to Hales' lemma 5.19 which allows us to apply the
descent argument. 
By continuity, we see that the factor
$\kappa(inv(\gamma_G^0,\gamma_{G,\lambda}))\kappa(t_*)$ satisfies the
local translation invariance property of lemma 5.19 above, for 
sufficiently small $Y$. One checks that the factor (8.7)
satisfies the translation invariance property using that
it is a product over roots in $H$. The theorem now follows from
lemma 5.19.

\enddocument



